% This work is licensed under the Creative Commons
% Attribution-NonCommercial-ShareAlike 4.0 International License. To view a copy
% of this license, visit http://creativecommons.org/licenses/by-nc-sa/4.0/ or
% send a letter to Creative Commons, PO Box 1866, Mountain View, CA 94042, USA.

% (c) Eric Kunze, 2020

%%%%%%%%%%%%%%%%%%%%%%%%%%%%%%%%%%%%%%%%%%%%%%%%%%%%%%%%%%%%%%%%%%%%%%%%%%%%%%%
% Template for exercises and homework at TU Dresden.
%%%%%%%%%%%%%%%%%%%%%%%%%%%%%%%%%%%%%%%%%%%%%%%%%%%%%%%%%%%%%%%%%%%%%%%%%%%%%%%

\documentclass[ngerman, a4paper, 11pt]{article}

\usepackage[top=2.5cm,bottom=2.5cm,left=2.5cm,right=2.5cm]{geometry}

\usepackage[ngerman]{babel}
\usepackage[utf8]{inputenc}

\usepackage{parskip}  	% split paragraphs by vspace instead of intendations
\usepackage[onehalfspacing]{setspace} % increase row-space

\usepackage{lmodern}
\usepackage{eufrak}
\usepackage[normalem]{ulem} 		% better underlines
\usepackage[autostyle=true,english=british]{csquotes}

\usepackage[scale=0.95]{opensans}	% new font OpenSans
%\newcommand*{\fosfamily}{\fontfamily{fos}\selectfont}
\DeclareTextFontCommand{\textos}{\fosfamily}

\usepackage{xifthen}


%%%%%%%%%%%%%%%%%%%%%%%%%%%%%%%%%%%%%%%%%%%%%%%%%%%%%%%%%%%%%%%%%%%%%%%%%%%%%%%%%%%%%%%%%%%%%%%%%%%%%%%

\newcommand{\name}{Eric Kunze}
\newcommand{\matnr}{4679202}
\newcommand{\email}{\href{mailto:eric.kunze@mailbox.tu-dresden.de}{\ttfamily eric.kunze@mailbox.tu-dresden.de}}

\newcommand{\modul}{Partielle Differentialgleichungen}
\newcommand{\semester}{Sommersemester 2020}

%\renewcommand{\tutor}{TUTOR}
%\renewcommand{\gruppe}{Tag x. DS, (un)gerade Woche}

\newcommand{\professor}{Prof. Dr. Friedemann Schuricht}
\newcommand{\fakultaet}{Mathematik}
\newcommand{\institut}{Geometrie}
\newcommand{\lehrstuhl}{Nichtlineare Analysis}

%%%%%%%%%%%%%%%%%%%%%%%%%%%%%%%%%%%%%%%%%%%%%%%%%%%%%%%%%%%%%%%%%%%%%%%%%%%%%%%%

\usepackage[smallequationskip]{../../texmf/tex/latex/mathworkMathTUD}
\usepackage{../../texmf/tex/latex/mathoperatorsMathTUD}

%%%%%%%%%%%%%%%%%%%%%%%%%%%%%%%%%%%%%%%%%%%%%%%%%%%%%%%%%%%%%%%%%%%%%%%%%%%%%%%%
% graphics

\usepackage{xcolor}
\usepackage[table,dvipsnames]{tudscrcolor}
\usepackage{graphicx}
\usepackage{tcolorbox}

\usepackage[font=small,labelfont=bf]{caption} % captions of non-floated figures

\usepackage{pgf}
\usepackage{pgfplots}
\pgfplotsset{compat=1.10}   % in my packages used compat=1.15
\usepgfplotslibrary{fillbetween}

\usepackage{tikz}
\usetikzlibrary{patterns,arrows,calc,decorations.pathmorphing,backgrounds, positioning,fit,petri,decorations.fractals}
\usetikzlibrary{matrix}


%%%%%%%%%%%%%%%%%%%%%%%%%%%%%%%%%%%%%%%%%%%%%%%%%%%%%%%%%%%%%%%%%%%%%%%%%%%%%%%%
% tabulars

\usepackage{tabularx} 
\usepackage{multirow}



%%%%%%%%%%%%%%%%%%%%%%%%%%%%%%%%%%%%%%%%%%%%%%%%%%%%%%%%%%%%%%%%%%%%%%%%%%%%%%%%
% counter

\usepackage{chngcntr}

% automatic reset of section after chapter ended 
\pretocmd{\chapter}{\setcounter{section}{0}}{}{}

% automatic reset of equation counter in each section
\pretocmd{\section}{\setcounter{equation}{0}}{}{}


\usepackage{chngcntr}
\newcounter{taskcount}
\newcounter{blattcount}
\newcounter{themcount}

\numberwithin{page}{blattcount}

\counterwithout{equation}{section}
\counterwithin{equation}{blattcount}
\counterwithin{figure}{blattcount}
\counterwithin{table}{blattcount}

%\counterwithin{themcount}{chapter}
%\numberwithin{page}{blattcount}
%\counterwithin{taskcount}{blattcount}

%%%%%%%%%%%%%%%%%%%%%%%%%%%%%%%%%%%%%%%%%%%%%%%%%%%%%%%%%%%%%%%%%%%%%%%%%%%%%%%%
% Environment for a new exercise sheet

\usepackage{tcolorbox}
\newcommand{\header}{%
        {\fosfamily%
            \fcolorbox{cddarkblue}{cdblue!20}{%
                \begin{minipage}{\dimexpr0.7\linewidth-\fboxrule-\fboxsep}
                    {\huge \textbf{Hausaufgaben}} \vspace{3pt} \\
                    \textbf{\modul} -- Übungsblatt \theblattcount
                \end{minipage}%
                \begin{minipage}{\dimexpr0.3\linewidth-\fboxrule-\fboxsep}
                    \flushright \textbf{\name} \\
                    Matr.-Nr. \matnr
                \end{minipage}%
        }}%
    }%

\NewDocumentEnvironment{exercisePage}{O{}}{%
    \pagebreak%
    \stepcounter{blattcount} \setcounter{page}{1}%
    \header \\[6pt]
    \setcounter{figure}{0}
    \setcounter{equation}{0}%
    \setcounter{table}{0}%
    \setcounter{themcount}{0}
    %
    \ifthenelse{\isempty{#1}}{}{%
        {\fosfamily \itshape Thema: #1} \par%\\[12pt]%
    }
}{}

%%%%%%%%%%%%%%%%%%%%%%%%%%%%%%%%%%%%%%%%%%%%%%%%%%%%%%%%%%%%%%%%%%%%%%%%%%%%%%%%
% Packages for theorem environments

\usepackage[ntheorem,framemethod=TikZ]{mdframed}
\usepackage{amsmath,amssymb,amsfonts,mathtools}
\usepackage[amsmath,thmmarks,framed]{ntheorem}

%%%%%%%%%%%%%%%%%%%%%%%%%%%%%%%%%%%%%%%%%%%%%%%%%
% theorem environments for exercises and solutions

\newcommand{\skiparound}{10pt}

% >> Exercises
\theoremstyle{plain}
\theoremheaderfont{\fosfamily\normalsize\bfseries\upshape}
\theorembodyfont{\normalsize}
\theoremseparator{.}
\theoremsymbol{}

\newmdtheoremenv[%
    backgroundcolor=cdblue!5,%
    linecolor=cddarkblue,%
    skipabove=\skiparound,%
    skipbelow=\skiparound,%
    nobreak,%
]{exercise}[taskcount]{Übung}

\newmdtheoremenv[%
    backgroundcolor=cdblue!5,%
    linecolor=cddarkblue,%
    skipabove=\skiparound,%
    skipbelow=\skiparound,%
    nobreak,%
]{homework}[taskcount]{Hausaufgabe}

\newmdtheoremenv[%
    backgroundcolor=cdblue!5,%
    linecolor=cddarkblue,%
    skipabove=\skiparound,%
    skipbelow=\skiparound,%
    nobreak,%
]{task}[taskcount]{Aufgabe}

\newmdtheoremenv[%
backgroundcolor=cdblue!5,%
linecolor=cddarkblue,%
skipabove=\skiparound,%
skipbelow=\skiparound,%
nobreak,%
]{zusatz}[taskcount]{Zusatzaufgabe}

%%%%%%%%%%%%%%%%%%%%%%%%%%%%%%%%%%%%%%%%%%%%
% >> Lemma
\theoremstyle{plain}
\theoremheaderfont{\fosfamily\normalsize\bfseries\upshape}
\theorembodyfont{}
\theoremseparator{.}
\theorempreskip{6pt}
\theorempostskip{6pt}
\newtheorem{lemma}{Lemma}
\counterwithin{lemma}{blattcount}

% >> Solution
\theoremstyle{nonumberplain}
\theoremheaderfont{\normalfont\normalsize\bfseries}
\theorembodyfont{\normalfont}
\theoremseparator{\textbf{.}}
\theorempreskip{5pt}
\theorempostskip{5pt}
\theoremsymbol{$\square$}
\newtheorem{solution}{Lösung}

% >> proof
\newtheoremstyle{proofstyle}%
{\item[\hskip\labelsep {\theorem@headerfont ##1}\theorem@separator]}%
{\item[\hskip\labelsep {\theorem@headerfont ##1}\ (##3)\theorem@separator]}

\theoremstyle{proofstyle}
\theoremheaderfont{\normalsize\slshape}
\theorembodyfont{}
\theoremseparator{.}
\theorempreskip{0pt}
\theorempostskip{5pt}
\theoremsymbol{$\square$}
\newtheorem{proof}{Beweis}

% >> Equivalences
\newcommand{\hinrichtung}{\item[\bfseries ($\boldsymbol{\Rightarrow}$)]}
\newcommand{\rueckrichtung}{\item[\bfseries ($\boldsymbol{\Leftarrow}$)]}

\newenvironment{equivalence}[1][]{%
    \ifthenelse{\isempty{#1}}%
        {\begin{description}}%
        {\begin{description}[topsep=-\parskip]}%
}{%
        \end{description}%
}%

% >> inductions
\newcommand{\ianfang}[1][]{%
    \ifthenelse{\isempty{#1}}{%
        % no parameter
        \item[\textbf{(IA)}]%
    }{%
        %parameter exists
        \item[\textbf{(IA)}] {#1 :} %\hfill \\%
    }%
}

\newcommand{\ivorraussetzung}{\item[\bfseries (IV)]}

\newcommand{\ischritt}[1][]{%
    \ifthenelse{\isempty{#1}}{%
        % no parameter
        \item[\textbf{(IS)}]
    }{%
        % parameter exists
        \item[\textbf{(IS)}] {#1 :} %\\
    }%
}

% two optional arg's:
% #1 = induced variable 
% #2 = vert. space before (IA): parskip (default) or nopreskip
\NewDocumentEnvironment{induction}{O{} O{}}{%
    \ifthenelse{ \isempty{#1} }{}{%
        Vollständige Induktion nach #1:
    }%---
    \ifthenelse{ \equal{#2}{nopreskip} }{%
        \begin{description}[topsep=-\parskip]
    }{%
        \begin{description}
    }%
}{%
    \end{description}%
}


%%%%%%%%%%%%%%%%%%%%%%%%%%%%%%%%%%%%%%%%%%%%%%%%%%%%%%%%%%%%%%%%%%%%%%%%%%%%%%%%
% titlepage

\newcommand{\makeTUtitle}[1][]{%
	\begin{titlepage}
		\pagecolor{cddarkblue!90} \color{white}%
		\raggedright \fosfamily%
		\setlength{\parindent}{0pt}% 
	% Logo / Kopf
		\hspace{-18.6mm} %
		\includegraphics[scale=0.6]{TUD-white.pdf} \\
		\vspace{3mm} 
		\begin{tabular}{m{\textwidth}}
			\hline
			\hspace{-4pt}\small{\textbf{Fakultät \fakultaet} Institut für \institut, Professur für \lehrstuhl} \\
			\hline
		\end{tabular} \\
	% Titel
		\vspace{5cm}
		{\Huge\bfseries \MakeUppercase \modul \par}
		\ifthenelse{\isempty{#1}}{}{
			\vspace{0.5cm}%
			{\Large \itshape #1} \\%
		}
		\vspace{1.5cm}
		\textbf{{\Large \professor}} \par
		\vspace{0.5cm}
		{\large \semester}
	% Fußzeile
		\vfill%
		\begin{tabular}{lll}
			Autor  & : & \name \\
			E-Mail & : & \email \\
		\end{tabular}%
	\end{titlepage}
	\nopagecolor
}

%%%%%%%%%%%%%%%%%%%%%%%%%%%%%%%%%%%%%%%%%%%%%%%%%%%%%%%%%%%%%%%%%%%%%%%%%%%%%%%%


%%%%%%%%%%%%%%%%%%%%%%%%%%%%%%%%%%%%%%%%%%%%%%%%%%%%%%%%%%%%%%%%%%%
%                          HIGHLIGHTING                           %
%%%%%%%%%%%%%%%%%%%%%%%%%%%%%%%%%%%%%%%%%%%%%%%%%%%%%%%%%%%%%%%%%%%
\newcommand{\begriff}[1]{\textbf{#1}}
\newcommand{\person}[1]{\textsc{#1}}

%%%%%%%%%%%%%%%%%%%%%%%%%%%%%%%%%%%%%%%%%%%%%%%%%%%%%%%%%%%%%%%%%%%
%                            EQUATIONS                            %
%%%%%%%%%%%%%%%%%%%%%%%%%%%%%%%%%%%%%%%%%%%%%%%%%%%%%%%%%%%%%%%%%%%

\setlength\abovedisplayshortskip{0pt plus 3pt}%
\setlength\belowdisplayshortskip{4pt plus 2pt minus 2pt}%
\setlength\abovedisplayskip{1pt plus 1pt minus 2pt}%
\setlength\belowdisplayskip{1pt plus 1pt minus 2pt}%


%%%%%%%%%%%%%%%%%%%%%%%%%%%%%%%%%%%%%%%%%%%%%%%%%%%%%%%%%%%%%%%%%%%
%                          ENUMERATIONS                           %
%%%%%%%%%%%%%%%%%%%%%%%%%%%%%%%%%%%%%%%%%%%%%%%%%%%%%%%%%%%%%%%%%%%
\usepackage{enumerate}
\usepackage[inline]{enumitem} 		%customize label

\renewcommand{\labelitemi}{\raisebox{1pt}{\scalebox{.5}{$\blacksquare$}}}
\renewcommand{\labelitemii}{$\vartriangleright$}
\renewcommand{\labelitemiii}{--}
% Variantionen des Dreiecks als Aufzählungszeichen $\blacktriangleright$ / $\vartriangleright$ / $\triangleright$

\renewcommand{\labelenumi}{(\alph{enumi})}
\renewcommand{\labelenumii}{\roman{enumii}.}
%\renewcommand{\labelenumiii}{\roman{enumiii}.}

%%%%%%%%%%%%%%%%%%%%%%%%%%%%%%%%%%%%%%%%%%%%%%%%%%%%%%%%%%%%%%%%%%%
%                            LISTINGS                             %
%%%%%%%%%%%%%%%%%%%%%%%%%%%%%%%%%%%%%%%%%%%%%%%%%%%%%%%%%%%%%%%%%%%
\usepackage{listings}

%%%%%%%%%%%%%%%%%%%%%%%%%%%%%%%%%%%%%%%%%%%%%%%%%%%%%%%%%%%%%%%%%%%%%%%%%%%%%%%%

%%%%%%%%%%%%%%%%%%%%%%%%%%%%%%%%%%%%%%%%%%%%%%%%%%%%%%%%%%%%%%%%%%%
%                           REFERENCES                            %
%%%%%%%%%%%%%%%%%%%%%%%%%%%%%%%%%%%%%%%%%%%%%%%%%%%%%%%%%%%%%%%%%%%

\usepackage[unicode,bookmarks=true]{hyperref}
\hypersetup{
	% pdfborder={0 0 0}			% no boxed around links
	pdfborderstyle={/S/U/W 1},	% underlining insteas of boxes
	linkbordercolor=cdblue,
	urlbordercolor=cdblue
	%	colorlinks,
	%	citecolor=black,
	%	filecolor=cddarkblue!80,
	%	linkcolor=black,
	%	urlcolor=cddarkblue!80
}


\usepackage{cleveref}
\crefname{lemma}{Lemma}{Lemmata}
\crefname{exercise}{Übung}{Übungen}
\crefname{homework}{Hausaufgabe}{Hausaufgaben}
\crefname{task}{Aufgabe}{Aufgaben}

\usepackage{bookmark}		% pdf-bookmarks

%%%%%%%%%%%%%%%%%%%%%%%%%%%%%%%%%%%%%%%%%%%%%%%%%%%%%%%%%%%%%%%%%%%%%%%%%%%%%%%%
%%%%%%%%%%%%%%%%%%%%%%%%%%%%%%%%%%%%%%%%%%%%%%%%%%%%%%%%%%%%%%%%%%%%%%%%%%%%%%%%
%%%%%%%%%%%%%%%%%%%%%%%%%%%%%%%%%%%%%%%%%%%%%%%%%%%%%%%%%%%%%%%%%%%%%%%%%%%%%%%%
%%%%%%%%%%%%%%%%%%%%%%%%%%%%%%%%%%%%%%%%%%%%%%%%%%%%%%%%%%%%%%%%%%%%%%%%%%%%%%%%
%%%%%%%%%%%%%%%%%%%%%%%%%%%%%%%%%%%%%%%%%%%%%%%%%%%%%%%%%%%%%%%%%%%%%%%%%%%%%%%%
%%%%%%%%%%%%%%%%%%%%%%%%%%%%%%%%%%%%%%%%%%%%%%%%%%%%%%%%%%%%%%%%%%%%%%%%%%%%%%%%


%---------------------------------------
% additional packages
%---------------------------------------

\usepackage{accents}

%---------------------------------------
% additional comands
%---------------------------------------

\renewcommand*{\dot}[1]{\accentset{\mbox{\large\bfseries .}}{#1}}
\renewcommand*{\ddot}[1]{\accentset{\mbox{\large\bfseries .\hspace{-0.25ex}.}}{#1}}

\newcommand{\dsq}[1]{\left[\left[ #1 \right]\right]}

%%%%%%%%%%%%%%%%%%%%%%%%%%%%%%%%%%%%%%%%%%%%%%%%%%%%%%%%%%%%%%%%%%%%%%%%%%%%%%%%%%%%%%%%%%%%%%%%%%%%%%%


\begin{document}

\makeTUtitle[Übungen]


\begin{exercisePage}
	
	\begin{task} Es sei $U \subset \Rn$ ein Gebiet, $u \in C^2(U,\Rn)$, $A \in C^1(U, \R^{m \times n})$, $\phi \in C^1(U)$ und $c \in \Rm$. Zeigen Sie:
		\begin{enumerate}
			\item $\div (\transpose{Du}) = \transpose{D(\div u)}$,
			\item $\div (cA) = c\div A$,
			\item $\div (\varphi A) = A * D\varphi + \varphi \div A$.
		\end{enumerate}
		\emph{Hinweis:} Wir betrachten Vektoren im $\Rn$ als Zeilenvektoren, $*$ ist das Skalarprodukt im $\Rn$, $Du = (\partial_j u_i)_{i,j=1,\dots,n}$, $\div$ und $*$ wirken auf eine Matrix \textit{zeilenweise}.
		
		Es seien nun $\abb{u,v}{U}{\R}$ hinreichend glatt. Beweisen Sie die \textbf{Formel von Leibniz}:
		\begin{equation*}
			D^\alpha(uv) = \sum_{\beta \leq \alpha} {\alpha \choose \beta} D^\beta u D^{\alpha-\beta}v
		\end{equation*}
		wobei für Multiindizes $\alpha, \beta$ gilt: 
		\begin{equation*}
			D^\alpha u = \partial_{x_1}^{\alpha_1}\dots\partial_{x_n}^{\alpha_n}u,\qquad {\alpha\choose\beta} = \frac{\alpha!}{\beta!(\alpha-\beta)!},\qquad \alpha! = \alpha_1!\cdot\alpha_2!\cdot\dots\cdots \alpha_n!
		\end{equation*}
		und $\beta \leq \alpha$ genau dann, wenn $\beta_i \leq \alpha_i$ für $i=1,\dots,n$.
	\end{task}

	\begin{enumerate}[label=(zu \alph*), leftmargin=*]
		\item Sei $u = (u_1, u_2, \dots, u_n) \in C^2$. Wir notieren Vektoren verkürzt $(u_i)_i = (u_i)_{i=1,\dots,n} = (u_1, u_2, \dots, u_n)$. Es ist
		\begin{align*}
			\div\brackets{\transpose{Du}} &= \div \trans{\begin{pmatrix}
				\partial_1 u_1 & \partial_2 u_1 & \cdots & \partial_n u_1 \\
				\vdots         & \vdots         & \ddots & \vdots \\
				\partial_1 u_n & \partial_2 u_n & \cdots & \partial_n u_n
			\end{pmatrix}} 
			= \div \transpose{\partial_j u_i}_{i,j}
			= \div \brackets{\partial_i u_j}_{i,j} \\
			&= \begin{pmatrix}
				\partial_{11} u_1 & \partial_{12} u_2 & \cdots & \partial_{1n} u_n \\
				\vdots         & \vdots         & \ddots & \vdots \\
				\partial_{n1} u_1 & \partial_{n2} u_2 & \cdots & \partial_{nn} u_n
			\end{pmatrix}
			= \brackets{\sum_{j=1}^n \partial_{ij} u_j}_{i}
		\end{align*}
		und außerdem
		\begin{equation*}
			\transpose{D(\div u)} 
			= \transpose{D \brackets{\sum_{i=1}^n \partial_i u_i}}
			= \begin{pmatrix}
				\partial_{11} u_1 + \partial_{21} u_2 + \dots + \partial_{n1} u_n \\
				\vdots \\
				\partial_{1n} u_1 + \partial_{2n} u_2 + \dots + \partial_{nn} u_n
			\end{pmatrix}
			= \brackets{\sum_{j=1}^n \partial_{ji} u_j}_{i}
		\end{equation*}
		Wegen $u \in C^2$ sind alle partiellen Ableitungen stetig und können somit vertauscht werden. Daraus folgt die (zeilenweise) Gleichheit mit
		\begin{equation*}
			\div\brackets{\transpose{Du}} = \brackets{\sum_{j=1}^n \partial_{ij} u_j}_{i} = \brackets{\sum_{j=1}^n \partial_{ji} u_j}_{i} = \transpose{D(\div u)} 
		\end{equation*} 
		%
		\item Es ist
		\begin{align*}
			\div(cA) &= \div \brackets{\sum_{j=1}^n c_j a_{ij}}_i = \sum_{i=1}^n \partial_i \sum_{j=1}^n c_j a_{ij} = \sum_{i=1}^n \sum_{j=1}^n \partial_i c_j a_{ij} \\
			c * \div(A) &= c * \brackets{\sum_{j=1}^n \partial_j a_{ij}}_i = \sum_{i=1}^n c_i \sum_{j=1}^n \partial_j a_{ij} = \sum_{i=1}^n \sum_{j=1}^n \partial_j c_j a_{ij}
		\end{align*}
		und somit $\div(cA) = c * \div(A)$.
		%
		\item Man hat
		\begin{align*}
			A * D\phi &= A * (\partial_i \phi)_i = \brackets{\sum_{j=1}^n a_{ij} \ \partial_j \phi}_i \\
			\phi \div(A) &= \phi * \brackets{\sum_{j=1}^n \partial_j a_{ij}}_i = \brackets{\sum_{j=1}^n \phi \ \partial_j a_{ij}}_i \\
			\div(\phi A) &= \div\brackets{\brackets{\phi a_{ij}}_{i,j}} = \brackets{\sum_{j=1}^n \partial_j (\phi a_{ij})}_i
		\end{align*}
		Für fixiertes $i$ (also zeilenweise) erhält man mit der Produktregel für (partielle) Ableitungen
		\begin{equation*}
			\sum_{j=1}^n \partial_j (\phi a_{ij}) = \sum_{j=1}^n \brackets{(\partial_j \phi) a_{ij} + \phi (\partial_j a_{ij})} = \sum_{j=1}^n (\partial_j \phi) a_{ij} + \sum_{j=1}^n \phi (\partial_j a_{ij})
		\end{equation*}
		und somit $\div(\phi A) = A * (D\phi) + \phi * \div(A)$.
	\end{enumerate}

	\textbf{Leibnitz-Formel:} Vollständige Induktion über $\abs{\alpha} = k$.
	\begin{induction}
		\ianfang[$k = 0$] Für $\abs{\alpha} = 0$, also $\alpha = 0$ ist 
		\begin{equation*}
			D^0(uv) = uv = \binom{0}{0} D^0 u \ D^0 v
		\end{equation*}
		\ivorraussetzung Für $\abs{\alpha} = \sum_{i=1}^n \alpha_i = k$ gilt $D^\alpha(uv) = \sum_{\beta \leq \alpha} {\alpha \choose \beta} D^\beta u D^{\alpha-\beta}v$.
		\pagebreak
		\ischritt[$k \to k+1$] Seien $\abs{\alpha} = \abs{(\alpha_1, \alpha_2, \dots, \alpha_n)} = k+1$ und $\alpha' = (\alpha_1, \dots, \alpha_{n-1}, \alpha_n -1)$ sowie $\beta' = (\beta_1, \dots, \beta_{n-1}, \beta_n + 1)$. Dann ist $\abs{\alpha'} = k$ und $\abs{\beta'} = \abs{\beta} + 1$. Es gilt
		\begin{align*}
			D^\alpha (uv) 
			&= \partial_{x_1}^{\alpha_1}\dots\partial_{x_n}^{\alpha_n} (uv)
			= \partial_{x_n} \brackets{\partial_{x_1}^{\alpha_1}\dots\partial_{x_n}^{\alpha_n - 1} (uv)} \\
			&= \partial_{x_n} \brackets{D^{\alpha'} (uv)} \\
			\overset{\text{IV}}&{=} \partial_{x_n} \brackets{\sum_{\beta \leq \alpha'} \binom{\alpha'}{\beta} D^\beta u \ D^{\alpha'-\beta}v} \\
			&= \sum_{\beta \leq \alpha'} \binom{\alpha'}{\beta} \brackets{\partial_{x_n} (D^\beta u) \ D^{\alpha'-\beta}v + D^\beta u \ \partial_{x_n} (D^{\alpha' - \beta} v) } \\
			&= \sum_{\beta \leq \alpha'} \binom{\alpha'}{\beta} \brackets{D^{\beta'} u \ D^{\alpha'-\beta}v + D^\beta u \ D^{\alpha - \beta} v } \\
			&= \dots \\
			&= \sum_{\beta \leq \alpha} {\alpha \choose \beta} D^\beta u \  D^{\alpha-\beta}v
		\end{align*}
	\end{induction}







	\begin{task}
		Ermitteln Sie jeweils eine nichttriviale Lösung der folgenden partiellen Differentialgleichungen:
		\begin{enumerate}
			\item $v_y(x,y)= xy * v(x,y)$
			\item $u_x(x,y) + y * u(x,y) = 0$
		\end{enumerate}
	\end{task}

	\begin{enumerate}[label=(zu \alph*), leftmargin=*]
		\item Wir gehen analog zur Vorlesung vor und betrachten die Gleichung für fixiertes $x = \text{const}$. Dann erhalten wir eine gewöhnliche Differentialgleichung
		\begin{equation*}
			u'(y) = x y * u(y) 
		\end{equation*} und lösen entweder durch geübtes Hinschauen oder mit Trennung der Variablen: sei $f(u(y)) = u(y)$ und $g(y) = x*y$. Der Ansatz
		\begin{equation*}
			\int^{u(y)} \frac{1}{f(\xi)} \diff{\xi} = \int^y g(\xi) \diff{\xi} \follows \ln(\abs{u(y)}) = \frac{1}{2} x y^2 + C
		\end{equation*} 
		Beachte, dass die unteren Integralgrenzen dabei in der Konstante $C$ zusammengefasst sind, da wir keinen Anfangswert vorgegeben haben. Die Gleichung \enquote{umgestellt} ergibt eine Lösung
		\begin{equation*}
			u(y) = \exp\brackets{\frac{1}{2}xy^2} \qquad \bzw \qquad v(x,y) = \exp\brackets{\frac{1}{2}xy^2}
		\end{equation*}
		Eine kurze Probe ergibt
		\begin{equation*}
			v_y(x,y) = xy * \exp\brackets{\frac{1}{2}xy^2} = xy * v(x,y)
		\end{equation*}
		%
		\item Wir fixieren erneut eine Variable, diesmal $y = \text{const}$. Diesmal sehen wir direkt eine Lösung, nämlich
		\begin{equation*}
			u(x,y) = \exp(-xy)
		\end{equation*}
		Eine kurze Probe ergibt $u_x(x,y) + y * u(x,y) = -y * \exp(-xy) + y * \exp(-xy) = 0$.
	\end{enumerate}

	\begin{zusatz}
		Klassifizieren Sie die nachstehenden partiellen Differentialgleichungen nach folgenden Gesichtspunkten:
		\begin{enumerate}[nolistsep, topsep=-\parskip]
			\item Ist die Differentialgleichung linear, semilinear, quasilinear oder voll nichtlinear?
			\item Welche Ordnung hat die Differentialgleichung?
		\end{enumerate}
		\begin{align*}
			\Delta u &=0\\
			-\Delta u &=f(u)\\
			|Du|&=1\\
			u_t + \sum_{i=1}^n b^i u_{x_i}&=0\\
			\det (D^2u) &=f\\
			\div \left( \frac{Du}{\sqrt{1+|Du|^2}} \right)&=0\\
			u_t-\Delta u &= f(u)\\
			u_{tt} - \Delta u&=0\\
			u_t + \div F(u)&=0 \\
			u_t + u u_x + u_{xxx} &=0 \\
			a u_t + \Delta u&=0\\
			u_t + H(Du,x)&=0\\
			u_t - \Delta (u^\gamma)&=0\\
			-\Delta u &= \lambda u \\
			\div(|Du|^{p-2}Du)&=0\\
			u_t + uu_x&=0
		\end{align*}
	\end{zusatz}
	
\end{exercisePage}
\begin{exercisePage}[Charakteristikenmethode]
	
	\begin{center} \itshape
		Ich bitte um Entschuldigung, dass meine Lösung so lang geworden ist. Aber ich habe mich bemüht mein Vorgehen detailliert zu beschreiben. \smiley
	\end{center}

	\begin{task}
		Finden Sie Lösungen $u\in C^1$ der folgenden linearen Randwertprobleme:
		\begin{enumerate}
			\item $-3u_x + 2u_y = 0$ mit $u(x,y) = y^2 + 1$ auf $\Gamma = \menge{(1,s) \in \R^2 ~|~ s \in \R}$
			\item $u_x + u_y - u_z = x e^{y-z}$ mit $u(0,y,z) = g(y,z)$ für alle $y,z \in \R$, wobei $g \in C^1(\R^2)$ beliebig ist
			\item $2u_x - u_y = 2u - x e^x$ mit $u(0,y) = y^2$ für alle $y \in \R$
		\end{enumerate}
		Erläutern Sie dabei Ihr Vorgehen und überprüfen Sie abschließend, ob die von Ihnen gefundene Lösung wirklich das Problem löst.
	\end{task}

	\begin{enumerate}[label=(zu \alph*), leftmargin=*]
		\item Wir betrachten die Charakteristiken
		\begin{equation*}
			\begin{pmatrix} x(t) \\ y(t) \end{pmatrix} = \begin{pmatrix} x_0 \\ y_0 \end{pmatrix} + t * \begin{pmatrix} -3 \\ 2 \end{pmatrix} \follows \begin{pmatrix} \dot{x}(t) \\ \dot{y}(t) \end{pmatrix} = \begin{pmatrix} -3 \\ 2 \end{pmatrix}
		\end{equation*}
		Sei $\alpha(t) \defeq u(x(t),y(t))$, d.h. $\alpha$ beschreibt $u$ entlang der Charakteristiken. Es gilt
		\begin{equation*}
			\alpha'(t) = u_x * \dot{x}(t) + u_y * \dot{y}(t) = -3u_x + 2u_y = 0
		\end{equation*}
		und somit ist $u$ konstant entlang der Charakteristiken. Parametrisiere die Kurve $\Gamma$ durch $x_0(s) = 1$ und $y_0(s) = s$. Dann ergibt sich die Randwertbedingung zu $g(s) = s^2 + 1$. Wir prüfen nun die nichtcharakteristische Bedingung, d.h. ob die Kurve $\Gamma$ auch alle Charakteristiken $\Xi_{(x_0,y_0)}$ durchläuft. Dazu prüfen wir den Tangentenvektor von $\Gamma$, nämlich $\left(\begin{smallmatrix} \dot{x}_0 \\ \dot{y}_0 \end{smallmatrix}\right) = \left(\begin{smallmatrix} 0 \\ 1 \end{smallmatrix}\right)$, und den Tangentenvektor der Charakteristik $\Xi_{(x_0,y_0)}$, nämlich $\left(\begin{smallmatrix} \dot{x} \\ \dot{y} \end{smallmatrix}\right) = \left(\begin{smallmatrix} -3 \\ 2 \end{smallmatrix}\right)$, auf lineare Unabhängigkeit:
		\begin{equation*}
			\det\begin{pmatrix} 0 & -3 \\ 1 & 2\end{pmatrix} = 3 \neq 0
		\end{equation*}
		Somit schneidet $\Gamma$ alle Charakteristiken $\Xi_{(x_0,y_0)}$. Somit können wir die Schar der Charakteristiken beschreiben durch 
		\begin{equation*}
			\begin{aligned}
				x(s,t) &= x_0(s) - 3t = 1 - 3t & \follows t(x,y) &= \frac{1-x}{3} \\
				y(s,t) &= y_0(s) + 2t = s + 2t & \follows s(x,y) &= y - 2t = y - \frac{2}{3} (1-x) = y + \frac{2}{3} x - \frac{2}{3}
			\end{aligned}
		\end{equation*}
		Nach Konstruktion in der Vorlesung erhalten wir damit eine Lösung 
		\begin{equation*}
			u(x,y) = g(s(x,y)) = \brackets{y + \frac{2}{3} x - \frac{2}{3}}^2 + 1
		\end{equation*}
		Die Probe liefert mit den partiellen Ableitungen
		\begin{equation*}
			\left. \begin{array}{rcrl}
			u_x(x,y) &=& \frac{4}{3} & \brackets{y + \frac{2}{3} x - \frac{2}{3}}^2 \\
			u_y(x,y) &=& 2  & \brackets{y + \frac{2}{3} x - \frac{2}{3}}^2
			\end{array} \right\} 
			\follows -3 * \frac{4}{3} \brackets{y + \frac{2}{3} x - \frac{2}{3}}^2 + 4 \brackets{y + \frac{2}{3} x - \frac{2}{3}}^2 = 0
		\end{equation*}
		und außerdem $u(1,y) = y^2 + 1$ für den Anfangswert.
		Somit ist $u$ also Lösung der Differentialgleichung.
		%
		\item Wir betrachten die partielle Differentialgleichung $u_x + u_y - u_z = x*e^{y-z}$ mit der Randbedingung $u(0,y,z) = g(y,z)$ für beliebiges $g \in C^1(\R^2)$. Definieren wir den \enquote{Rand} als die Fläche $\Gamma = \menge{(0,y,z) \in \R^3 : y,z \in \R}$. Diese lässt sich parametrisieren mit
		\begin{equation*}
			\gamma(\sigma, \tau) = 
			\begin{pmatrix} 
				x_0(\sigma, \tau) \\ 
				y_0(\sigma, \tau) \\ 
				z_0(\sigma, \tau) 
			\end{pmatrix} 
			= 
			\begin{pmatrix}
				0 \\
				\sigma \\
				\tau
			\end{pmatrix}
		\end{equation*}
		Für die Tangentialebene erhalten wir die Spannvektoren 
		\begin{equation*}
			\gamma_\sigma(\sigma, \tau) = 
			\begin{pmatrix}
				0 \\ 1 \\ 0
			\end{pmatrix}
			\qquad
			\gamma_\tau(\sigma, \tau) = 
			\begin{pmatrix}
				0 \\ 0 \\ 1
			\end{pmatrix}
			\follows
			\dot{\gamma}(\sigma, \tau) = 
			\begin{pmatrix}
				0 & 0 \\ 1 & 0 \\ 0 & 1
			\end{pmatrix}
		\end{equation*}
		Betrachten wir die Charakteristiken $\Xi_{\sigma, \tau} = \Image(\xi)$ mit
		\begin{equation} \label{2eq: charakteristiken}
			\xi(t, \sigma, \tau) =
			\begin{pmatrix}
				x(t, \sigma, \tau) \\
				y(t, \sigma, \tau) \\
				z(t, \sigma, \tau)
			\end{pmatrix}
			= 
			\begin{pmatrix}
				x_0(\sigma, \tau) + t \\
				y_0(\sigma, \tau) + t \\
				z_0(\sigma, \tau) - t \\
			\end{pmatrix}
			=
			\begin{pmatrix}
				t \\
				\sigma + t \\
				\tau - t \\
			\end{pmatrix}
		\end{equation}
		Prüfen wir die nichtcharakteristische Bedingung um sicherzustellen, dass auch jede Charakteristik $\Xi_{\sigma, \tau}$ von $\Gamma$ durchlaufen wird:
		\begin{equation*}
			\det(\dot{\gamma} ~|~ \dot{\xi}) = \det\begin{pmatrix}
				0 & 0 & 1 \\ 1 & 0 & 1 \\ 0 & 1 & -1
			\end{pmatrix}
			= 1 \neq 0
		\end{equation*}
		Bezeichne mit $f(u,x,y,z) = x * e^{y-z}$ die rechte Seite der Differentialgleichung. Schreibe $\xi(t) = \xi(t, \sigma, \tau)$ für fixiertes $\sigma$ und $\tau$ ($x$,$y$,$z$ analog). Sei $\alpha(t) \defeq u(\xi(t))$ die Funktion $u$ entlang einer Charakteristik $\Xi_{\sigma, \tau}$. Dann gilt
		\begin{equation} \label{2eq: ode}
			\begin{aligned}
				\dot{\alpha}(t) 
				&= Du * \dot{\xi} = f(u(\xi(t)), \xi(t)) = f(\alpha(t), t, \sigma + t, \tau -t) = t * e^{(\sigma + t) - (\tau - t)} \\
				&= t*e^{\sigma - \tau + 2t}
			\end{aligned}
		\end{equation}
		und dem Anfangswert $\alpha(0) = \alpha(0, \sigma, \tau) = g(\sigma, \tau)$. Lösen wir also dieses Anfangswertproblem und integrieren dazu die rechte Seite in \cref{2eq: ode} partiell:
		\begin{equation*}
			\begin{aligned}
				\alpha(t) 
				&= \int t*e^{\sigma - \tau + 2t} \diff{t} 
				= \brackets{\frac{1}{2} e^{\sigma - \tau + 2t}} t - \int \frac{1}{2} e^{\sigma - \tau + 2t} \diff{t} \\
				&= \frac{1}{2} t * e^{\sigma - \tau + 2t} - \frac{1}{4} e^{\sigma - \tau + 2t} + C(\sigma, \tau) \\
				&= \brackets{\frac{1}{2} t - \frac{1}{4}} e^{\sigma - \tau + 2t} + C(\sigma, \tau)
			\end{aligned}
		\end{equation*}
		Mit dem Anfangswert $\alpha(0) = g(\sigma, \tau)$ ergibt sich die Konstante
		\begin{equation*}
			\alpha(0) = \frac{1}{4} e^{\sigma - \tau} + C(\sigma, \tau) \overset{!}{=} g(\sigma, \tau) 
			\follows
			C(\sigma, \tau) = \frac{1}{4} e^{\sigma - \tau} + g(\sigma, \tau)
		\end{equation*}
		und somit die konkrete Lösung
		\begin{equation*}
			\alpha(t) = \brackets{\frac{1}{2} t - \frac{1}{4}} e^{\sigma - \tau + 2t} + \frac{1}{4} e^{\sigma - \tau} + g(\sigma, \tau)
		\end{equation*}
		Aus \cref{2eq: charakteristiken} erhalten wir die Inverse von $\xi$ als
		\begin{equation*}
			\xi^{-1}(x,y,z) =
			\begin{pmatrix}
				t(x,y,z) \\ \sigma(x,y,z) \\ \tau(x,y,u) 
			\end{pmatrix}
			=
			\begin{pmatrix}
				x \\ y-x \\ z+x
			\end{pmatrix}
		\end{equation*}
		Nach Konstruktion in der Vorlesung erhalten wir die Lösung
		\begin{equation*}
			\begin{aligned}
				u(x,y,z) 
				&= \alpha(\xi^{-1}(x,y,z)) \\
				&= \brackets{\frac{1}{2} x - \frac{1}{4}} e^{(y-x) - (z+x) + 2x} + \frac{1}{4} e^{(y-x)-(z+x)} + g(y-x, z+x) \\
				&= \brackets{\frac{1}{2} x - \frac{1}{4}} e^{y-z} + \frac{1}{4} e^{y- z -2x} + g(y-x, z+x) 
			\end{aligned} 
		\end{equation*}
		Nun prüfen wir noch, dass die gefundene Funktion auch wirklich eine Lösung der Differentialgleichung ist. Für die partiellen Ableitungen gilt
		\begin{equation*}
			\begin{aligned}
				u_x(x,y,z) &= \frac{1}{2}e^{y-z} - \frac{1}{2} e^{y-z-2x}  - \partial_1 g(y-x, z+x) + \partial_2 g(y-x, z+x) \\
				u_y(x,y,z) &= \brackets{\frac{1}{2}x - \frac{1}{4}}e^{y-z} + \frac{1}{4} e^{y-z-2x} + \partial_1 g(y-x, z+x) \\
				u_z(x,y,z) &= - \brackets{\frac{1}{2}x - \frac{1}{4}}e^{y-z} - \frac{1}{4} e^{y-z-2x} + \partial_2 g(y-x, z+x)
			\end{aligned}
		\end{equation*}
		Einsetzen liefert
		\begin{equation*}
			\begin{aligned}
				& \frac{1}{2}e^{y-z} - \frac{1}{2} e^{y-z-2x}  - \partial_1 g(y-x, z+x) + \partial_2 g(y-x, z+x) \\
				+ \ &
				\brackets{\frac{1}{2}x - \frac{1}{4}}e^{y-z} + \frac{1}{4} e^{y-z-2x} + \partial_1 g(y-x, z+x) \\
				+ \ &
				\brackets{\frac{1}{2}x + \frac{1}{4}}e^{y-z} + \frac{1}{4} e^{y-z-2x} - \partial_2 g(y-x, z+x) \\
				= \ &
				e^{y-z} * \brackets{\frac{1}{2} + \frac{1}{2}x - \frac{1}{4} + \frac{1}{2}x - \frac{1}{4}}  - e^{y-z-2x} \brackets{\frac{1}{2} - \frac{1}{4} - \frac{1}{4} } \\
				= \ &
				x * e^{y-z}
			\end{aligned}
		\end{equation*}
		Außerdem ist die Randwertbedingung erfüllt, denn
		\begin{equation*}
			u(0,y,z) = -\frac{1}{4} e^{y-z} + \frac{1}{4} e^{y-z} + g(y, z) = g(y,z)
		\end{equation*}
		und somit $u$ tatsächlich Lösung der partiellen Differentialgleichung.
		%
		\item Gegeben sei die partielle Differentialgleichung $2u_x - u_y = 2u + x*e^x$ und die Randwertbedingung $u(0,y) = y^2$ für alle $y \in \R$. Bezeichnen wir mit $f(u,x,y) = 2u - x e^x$  die rechte Seite.  Die Randwerte werden auf der Kurve $\Gamma = \menge{(0,y) \in \R^2 : y \in \R}$ angenommen. Diese können wir parametrisieren mit 
		\begin{equation*}
			\gamma(s) = 
			\begin{pmatrix} 
				x_0(s) \\ 
				y_0(s) 
			\end{pmatrix} 
			= 
			\begin{pmatrix}
				0 \\
				s
			\end{pmatrix}
			\follows \dot{\gamma}(s) =
			\begin{pmatrix} 
				\dot{x}_0(s) \\ 
				\dot{y}_0(s) 
			\end{pmatrix} 
			= 
			\begin{pmatrix}
				0 \\
				1
			\end{pmatrix}
		\end{equation*}
		Damit wird die Randwertbedingung zu $g(s) = s^2$. Betrachten wir die Charakteristiken $\Xi_s$ mit
		\begin{equation} \label{2eq: c-charakteristiken}
			\xi(t, s) =
			\begin{pmatrix}
				x(t, s) \\
				y(t, s)
			\end{pmatrix}
			= 
			\begin{pmatrix}
				x_0(s) + 2t \\
				y_0(s) - t \\
			\end{pmatrix}
			=
			\begin{pmatrix}
				2t \\
				s + t \\
			\end{pmatrix}
		\end{equation}
		Die nichtcharakteristische Bedingung ist hier erfüllt, denn
		\begin{equation*}
			\det \begin{pmatrix}
				0 & 2 \\ 1 & -1
			\end{pmatrix}
			= -2 \neq 0
		\end{equation*}
		Betrachte nun die Funktion $u$ entlang der Charakteristiken $\Xi_s$ für fixiertes $s$ beschrieben durch $\alpha(t) = u(\xi(t))$. Differenzieren ergibt
		\begin{equation} \label{2eq: c-anfangswertproblem}
			\dot{\alpha}(t) = u_x * \dot{x} + u_y * \dot{y} = f(u(\xi(t)), \xi(t)) = f(\alpha(t), 2t, s+t) = 2\alpha - 2t*e^{2t}
		\end{equation} 
		bei $\alpha(0,s) = g(s) = s^2$. Dieses Anfangswertproblem lösen wir mit Variation der Konstanten. Das zugehörige homogene Problem besitzt offensichtlich die Lösung $\alpha(t) = c(t) * e^{2t}$. Differenzieren wir diese Gleichung erhalten wir $\dot{\alpha}(t) = \dot{c}(t) * e^{2t} + 2c(t) * e^{2t}$. Setzen wir dies nun in \cref{2eq: c-anfangswertproblem} ein, dann erhalten wir für ein $\hat{c} \in \R$
		\begin{equation*}
			\dot{c}(t) * e^{2t} + 2c(t) * e^{2t} = 2c(t) * e^{2t} - 2t*e^{2t} \follows \dot{c}(t) = -2t \follows c(t) = \hat{c} - t^2
		\end{equation*}
		Damit ergibt sich die allgemeine Lösung $\alpha(t) = e^{2t} (\hat{c} - t^2)$. Durch den Anfangswert gilt $\alpha(0) = \hat{c} = s^2$ und somit ist $\alpha(t) = e^{2t} ( s^2 - t^2)$ konkrete Lösung des Anfangswertproblems, was sich auch leicht überprüfen lässt:
		\begin{equation*}
			\dot{\alpha}(t) = 2\underbrace{e^{2t}(s^2 - t^2)}_{= \alpha(t)} - 2t * e^{2t} = 2\alpha(t) - 2t * e^{2t} \quad \und \quad \alpha(0) = s^2
		\end{equation*}
		Aus \cref{2eq: c-charakteristiken} erhalten wir 
		\begin{equation*}
			\xi^{-1}(x,y) = 
			\begin{pmatrix}
				t(x,y) \\ s(x,y)
			\end{pmatrix}
			=
			\begin{pmatrix}
				\frac{1}{2} x \\ \frac{1}{2} x + y
			\end{pmatrix}
		\end{equation*}
		Damit folgt nach Konstruktion in der Vorlesung eine Lösung
		\begin{equation*}
			\begin{aligned}
				u(x,y) = \alpha(s(x,y)) 
				&= e^{x} \brackets{\brackets{\frac{1}{2} x + y}^2 - \frac{1}{4} x^2} \\
				&= e^x \brackets{\frac{1}{4} x^2 + xy + y^2 - \frac{1}{4} x^2 } \\
				&= e^x \brackets{y^2 + xy} 
			\end{aligned}
		\end{equation*}
		Dann gilt für die partiellen Ableitungen
		\begin{equation*}
			\begin{aligned}
				u_x &= e^x \brackets{y^2 + xy} + y*e^x \\
				u_y &= e^x \brackets{2y + x} 
			\end{aligned}
		\end{equation*}
		und somit
		\begin{equation*}
			\begin{aligned}
				2u_x - u_y &= 2e^x (y^2 + xy) + 2y * e^x - e^x (2y + x) \\
				&= 2u + e^x (2y - 2y - x) \\
				&= 2u - x*e^x 
			\end{aligned}
		\end{equation*}
		und $u(0,y) = e^0 (y^2 + 0*y) = y^2$. Damit ist also $u$ tatsächlich Lösung der partiellen Differentialgleichung.
	\end{enumerate}
\end{exercisePage}
\begin{exercisePage}
	\stepcounter{taskcount}
	\begin{task}
		Bestimmen Sie eine Lösung $u \in C¹(U)$ des quasilinearen Randwertproblems
		\begin{equation*}
			\begin{aligned}
				u u_x + u_y &= 1 \\
				u(x,x) &= \frac{1}{2}x \quad \forall x \in \R\setminus\menge{\xi} 
			\end{aligned} 
		\end{equation*}
		wobei $\xi \in \R$ geeignet gewählt und $U$ eine geeignet gewählte Umgebung der Menge ist, auf der $u$ vorgegeben ist. Nutzen Sie dazu die Methode der Charakteristiken, überprüfen Sie Ihr Ergebnis und skizzieren Sie einige Charakteristiken in der Nähe des Punktes $(\xi,\xi)$.
	\end{task}

	Aus der Vorlesung kennen wir die Notation $	a(u(x),x) * D u + b(u(x),x) = 0$. Wir notieren
	$a(u(x,y),x,y) = (u(x,y), 1)$ und $b(u(x,y), x, y) = -1$.   Aus der Randwertbedingung erhalten wir eine Kurve $\Gamma = \menge{(x,x) \in \R^2 : x \in \R\setminus\menge{\xi}}$ mit Parametrisierung $\gamma(s) = \left(\begin{smallmatrix} x_0(s) \\ y_0(s) \end{smallmatrix}\right) = \left(\begin{smallmatrix} s \\ s \end{smallmatrix}\right)$, auf der $g(s) = \frac{1}{2} s$ gilt. Wir überprüfen die nichtcharakteristische Bedingung gemäß Konstruktion in der Vorlesung als
	\begin{equation*}
		\det \brackets{ \dot{\gamma} \mid a(g(s),\gamma(s))} = \det\begin{pmatrix} 1 & g(s) \\ 1 & 1 \end{pmatrix} = \det\begin{pmatrix} 1 & \frac{1}{2}s \\ 1 & 1 \end{pmatrix} = 1 - \frac{1}{2} s \neq 0 \enskip \forall s \neq 2
	\end{equation*}
	Wähle somit also $\xi = 2$, um die Regularität zu sichern. Betrachten wir $\alpha(t,s) = u(x(t,s), y(t,s)$ als die Funktion $u$ entlang der Charakteristiken. Da die partielle Differentialgleichung quasilinear ist, reichen die beiden folgenden charakteristischen Gleichungen zu lösen aus:
	\begin{equation*}
		\begin{aligned}
			\begin{pmatrix} \dot{x} \\ \dot{y} \end{pmatrix} &= a ( \alpha, x,y) = \begin{pmatrix} \alpha \\ 1 \end{pmatrix} \\
			\dot{\alpha} = - b(\alpha, x, y)  = 1
		\end{aligned}
	\end{equation*}
	mit den Anfangswerten $\alpha(0,s) = \frac{1}{2} s$, $x(0,s) = s$ und $y(0,s) = s$. Lösen wir diese gewöhnlichen Differentialgleichungen:
	\begin{equation*}
		\setlength{\arraycolsep}{1pt}
		\begin{array}{rclcrclcrcl}
			\alpha(t) &=& t + c(s) &\quad\overset{\text{AW}}{\follows}\quad& c(s) &= &\frac{1}{2} s &\quad\follows\quad& \alpha(t,s) &=& t + \frac{1}{2} s \\
			y(t) &=& t + c(s) &\quad\overset{\text{AW}}{\follows}\quad& c(s) &=& s &\quad\follows\quad& y(t,s) &=& t + s \\
			x(t) &=& \frac{1}{2} t^2 + \frac{1}{2} s t + c(s) &\quad\overset{\text{AW}}{\follows}\quad& c(s) &=& s &\quad\follows\quad& x(t,s) &=&  \frac{1}{2} t^2 + \frac{1}{2} s t + s
		\end{array}
	\end{equation*}
	Wegen der charakteristischen Bedingung können wir dieses Gleichungssystem nach $t$ und $s$ auflösen:
	\begin{equation*}
		\begin{aligned}
			y(t,s) &= t + s \follows s = y - t \\
			x(t,s) &= \frac{1}{2}t^2 + \frac{1}{2} s t + s = \frac{1}{2}t^2 + \frac{1}{2} (y - t) t + y - t = t \brackets{\frac{1}{2} y - 1} + y
		\end{aligned}
	\end{equation*}
	also
	\begin{equation*}
		t(x,y) = \frac{x-y}{\frac{1}{2} y - 1} \quad \und \quad s(x,y) = y - \frac{x-y}{\frac{1}{2} y - 1}
	\end{equation*}
	Setzen wir dies als Lösung $u(x,y) = \alpha(t(x,y),s(x,y)$ zusammen, erhalten wir
	\begin{equation*}
		\begin{aligned}
			u(x,y) &= \alpha\brackets{\frac{x-y}{\frac{1}{2}y - 1} \ , \ y - \frac{x-y}{\frac{1}{2} y - 1}} \\
			&= \frac{x-y}{\frac{1}{2}y - 1} + \frac{1}{2} y - \frac{1}{2} \frac{x-y}{\frac{1}{2} y - 1} \\
			&= \frac{x-y}{y - 2} + \frac{1}{2} y \qquad \forall x, y \in \R, y \neq 2
		\end{aligned}
	\end{equation*} 
	Überprüfen wir unser Ergebnis:
	Es gilt
	\begin{equation*}
		\begin{aligned}
			u(x,y) &= \frac{x-y}{y - 2} + \frac{1}{2} y \\
			u_x(x,y) &= \frac{1}{y-2} \\
			u_y(x,y) &= \frac{-1}{y-2} -  \frac{x-y}{(y-2)^2} + \frac{1}{2} =\frac{2-x}{(y-2)^2} + \frac{1}{2}
		\end{aligned}
	\end{equation*}
	Setzen wir dies in die partielle Differentialgleichung ein, so erhalten wir
	\begin{equation*}
		\begin{aligned}
			u(x,y) * u_x(x,y) + u_y(x,y) &= \brackets{\frac{x-y}{y - 2} + \frac{1}{2} y} * \frac{1}{y-2} + \frac{2-x-2y}{(y-2)^2} + \frac{1}{2} \\
			&= \frac{x-y}{(y-2)^2} + \frac{1}{2} y \frac{1}{y-2} + \frac{2-x}{(y-2)^2} + \frac{1}{2} \\
			&= \frac{2-y}{(y-2)^2} + \frac{1}{2} y \frac{1}{y-2} + \frac{1}{2} \\
			&= \brackets{\frac{1}{2} y - 1} \frac{1}{y-2} + \frac{1}{2} \\
			&= 1
		\end{aligned}
	\end{equation*}
	und für die Randwerte $u(x,x) = \frac{1}{2} x$ für alle $x \neq 2$.
	
	Betrachten wir die Charakteristiken beschrieben mit einer Parametrisierung für fixiertes $s \neq 2$ und betrachten das Gleichungssystem
	\begin{equation*}
		\begin{aligned}
			x(t) &= \begin{pmatrix} \frac{1}{2} t^2 + \frac{1}{2} s t + s \\ t + s \end{pmatrix} 
			= \begin{pmatrix} 2 \\ 2 \end{pmatrix} 
			\follows t &= 2 - s
			\follows 2 &= \frac{1}{2} (2-s)^2 + \frac{1}{2} s (2-s) + s = 2
		\end{aligned}
	\end{equation*}
	Somit sind die Gleichungen unabhängig von $s \in \R \setminus \menge{2}$ erfüllt und alle Charakteristiken gehen durch den Punkt $(2,2)$.  
	
	\pgfplotsset{
		  compat=1.10,% mit writeLaTeX bisher noch nicht möglich
		  every axis/.append style={
		  	axis x line=middle,    % put the x axis in the middle
		  	axis y line=middle,    % put the y axis in the middle
		  	axis line style={->}, % arrows on the axis
		  	enlargelimits=0.05,
		  	grid,
		  	domain=-5:5,
		  }
		}
		
	\centering
	\begin{tikzpicture} 
		\begin{axis}[%
			xmin =   0,
			xmax =  4,
			ymin  = -2,
			ymax =  4,
			axis equal,
			legend style={
				at={(0,0)},
				anchor=north west,at={(axis description cs:1.1,1)}}
		]
		% s = 0
		\addplot [domain=-4:4, smooth, samples=100, variable=\t, cddarkblue, line width=1pt]({0.5*t^2}, {t});
		% s = 0.5
		\addplot [domain=-4:4, smooth, samples=100, variable=\t, cddarkgreen, line width=1pt]({0.5*t^2 + 0.5 * t * 0.5 + 0.5}, {t + 0.5});
		% s = 1
		\addplot [domain=-4:4, smooth, samples=100, variable=\t, cdpurple, line width=1pt]({0.5*t^2 + 0.5 * t * 1 + 1}, {t + 1});
		\addplot [domain=-5:4, smooth, samples=100, variable=\t, cdorange, line width=1pt]({0.5*t^2 + 0.5 * t * 3 + 3}, {t + 3});
		\addplot [domain=-5:4, smooth, samples=100, variable=\t, cdblue, line width=1pt]({0.5*t^2 + 0.5 * t * 4 + 4}, {t + 4});
		\legend{$s=0$,$s=\frac{1}{2}$,$s=1$,$s=3$,$s=4$}
		\end{axis} 
		\node[above,font=\bfseries] at (current bounding box.north) {Charakteristiken $x(t) = \begin{pmatrix} \frac{1}{2} t^2 + \frac{1}{2} s t + s \\ t + s \end{pmatrix}$};
	\end{tikzpicture} 
	
\end{exercisePage}
\begin{exercisePage}
\begin{task}
	Wir betrachten die \textit{Burgers-Gleichung}
	\begin{equation*}
	\begin{aligned}
	u_t + uu_x &= 0 \qquad \text{für } (x,t) \in \R \times (0,\infty)\\
	u(x,0) &= g(x) \qquad\text{für alle } x \in \R
	\end{aligned}
	\end{equation*}
	\begin{enumerate}
		\item Es sei zuerst $g$ gegeben durch
		\begin{equation*}
			g(x) = \begin{cases} 1 &\text{für } x \geq 0,\\ 0 & \text{für } x < 0 \end{cases}
		\end{equation*}
		Zeigen Sie, dass die Funktionen
		\begin{equation*}
			v(x,t) = \begin{cases} 1 & \text{für } x \geq \tfrac12 t,\\ 0 &\text{für }x < \tfrac12t \end{cases}
			\qquad \und \qquad
			w(x,t) = \begin{cases} 1 & \text{für } x \geq t,\\ \frac{x}{t} & \text{für } x\in [0,t),\\ 0 & \text{für } x < 0 \end{cases}
		\end{equation*}
		schwache Lösungen sind.
		
		\textit{Hinweis}: Es genügt, die Randwerte, die Rankine-Hugoniot-Bedingung auf Sprungkurven und die Differentialgleichung abseits dieser Kurven zu überprüfen.
		
		\item Es sei nun $g$ gegeben durch
		\begin{equation*}
			g(x) = \begin{cases} 1 & \text{für } x\leq 0 \\ 1-x & \text{für } 0 \leq x \leq 1 \\ 0 & \text{für } x \geq 1 \end{cases}
		\end{equation*}
		\begin{enumerate}
			\item Offenbar ist $g \notin C^1(\R)$. Stellen Sie dennoch die charakteristischen Gleichungen auf und ermitteln Sie, welche Lösung man formal für $t<1$ erwarten würde.
			\item In welchen Punkten $(x,t)\in \R \times (0,1)$ gilt $u_t + u u_x = 0$?
			\item Setzen Sie $u$ derart auf $\R \times (0,\infty)$ fort, dass nur eine Sprungkurve existiert und dass entlang dieser die Rankine-Hugoniot-Bedingung erfüllt ist.
			
			\textit{Hinweis}: Es bietet sich an, dass $u$ für $t \geq 1$ nur die Werte $0$ und $1$ annimmt.
		\end{enumerate}
	\end{enumerate}
\end{task}

Für die Burgers-Gleichung gilt in Anlehnung an die Notation der Vorlesung $F(u)_x = u u_x$ und damit $F(u) = \frac{1}{2} u^2$.
\begin{enumerate}[label=(zu \alph*), leftmargin=*]
	\item Wir betrachten die Funktion $v$. Für die Randwerte, d.h. für $y=0$ gilt 
	\begin{equation*}
		v(x,0) = \begin{cases} 1 &\text{für } x \geq 0 \\ 0 & \text{für } x < 0 \end{cases} \enskip = g(x) \qquad \text{für alle } x \in \R
	\end{equation*}
	Die Sprungkurve ist gegeben durch $s(t) = \frac{1}{2} t$ mit $\dot{s} = \frac{1}{2}$. Es gilt $\dsq{v} = -1$ und $\dsq{F(v)} = 0 - \frac{1}{2} = \frac{1}{2}$, also ist mit $\dsq{F(v)} = \dot{s} * \dsq{v}$ die Rankine-Hugoniot-Bedingung erfüllt.
	Betrachte die Differentialgleichung abseits der Sprungkurve:
	\begin{itemize}
		\item Sei $x > \frac{1}{2} t$. Dann ist $v$ gegeben durch $v(x,t) = 1$ und somit $v_t = v_x \equiv 0$. Eingesetzt in die PDE ergibt dies $v_t + v v_x = 0 + 1 * 0 = 0$. $\checkmark$ 
		\item Sei $x < \frac{1}{2} t$. Dann ist $v$ gegeben durch $v(x,t) = 0$ und somit $v_t = v_x \equiv 0$. Eingesetzt in die PDE ergibt dies $v_t + v v_x = 0 + 0 * 0 = 0$. $\checkmark$ 
	\end{itemize}

	Betrachten wir nun die Funktion $w$. Die Randwerte werde wegen
	\begin{equation*}
		w(x,0) = \begin{cases} 1 &\text{für } x \geq 0 \\ 0 & \text{für } x < 0 \end{cases} \enskip = g(x) \qquad \text{für alle } x \in \R
	\end{equation*}
	erfüllt. Wir erhalten hier zwei Sprungkurven:
	\begin{itemize}
		\item Die Winkelhalbierende können wir durch $s_1(t) = t$ parametrisieren, also ist $\dot{s} = 1$. Dann ergibt sich für die Differenzen
		\begin{equation*}
			\dsq{F(w)} = \frac{1}{2} - \frac{1}{2} = 0  \quad \und \quad \dsq{w} = 1 - \frac{t}{t} = 0 \quad \follows \quad \dsq{F(w)} = \dot{s_1} * \dsq{w}
		\end{equation*}
		\item Für die andere Sprungkurve, die wir mit $s_2(t) = 0$ parametrisieren ($\dot{s} = 0$), erhalten wir 
		\begin{equation*}
		\dsq{F(w)} = 0 - 0 = 0  \quad \und \quad \dsq{w} = \frac{0}{t} - 0 = 0 \quad \follows \quad \dsq{F(w)} = \dot{s_2} * \dsq{w}
		\end{equation*}
	\end{itemize}
	Die Differentialgleichung wird abseits der Sprungkurven auch erfüllt:
	\begin{itemize}
		\item Sei $x > t$. Dann ist $w$ gegeben durch $w(x,t) = 1$ und somit $w_t = w_x \equiv 0$. Eingesetzt in die PDE ergibt dies $w_t + w w_x = 0 + 1 * 0 = 0$. $\checkmark$ 
		\item Sei $x \in [0,t)$. Dann ist $w$ gegeben durch $w(x,t) = \frac{x}{t}$ und somit $w_t = \frac{- x}{t^2}$ und $w_x = \frac{1}{t}$. Eingesetzt in die PDE ergibt dies $w_t + w w_x = \frac{- x}{t^2} + \frac{x}{t} \frac{1}{t} = 0$. $\checkmark$ 
		\item Sei $x < 0$. Dann ist $w$ gegeben durch $w(x,t) = 0$ und somit $w_t = w_x \equiv 0$. Eingesetzt in die PDE ergibt dies $w_t + w w_x = 0 + 0 * 0 = 0$. $\checkmark$ 
	\end{itemize}

	Damit sind also $v$ und $w$ schwache Lösungen der partiellen Differentialgleichung.
	
	\item Mit $y \defeq (x,t)$, $a(u,y) = \left( \begin{smallmatrix} u \\ 1 \end{smallmatrix} \right)$ und $b(u,y) = 0$ hat die Burgergsgleichung die quasilineare Form der Vorlesung. Dann sind die charakteristischen Gleichungen gegeben durch
	\begin{equation*}
		\begin{aligned}
			\dot{y}(\tau, \sigma) &= \left( \begin{smallmatrix} \dot{x} \\ \dot{t} \end{smallmatrix} \right) = a(\alpha, y) = \left( \begin{smallmatrix} \alpha \\ 1 \end{smallmatrix} \right) \\
			\dot{\alpha}(\tau, \sigma) &= a(\alpha, y) * Du = - b(u,y) = 0 \enskip \mit \enskip \alpha(0,\sigma) = g(\sigma)
		\end{aligned}
	\end{equation*}
	Aus der zweiten Gleichung erhalten wir die Lösung $\alpha(\tau, \sigma) = g(\sigma)$, d.h. $u$ ist entlang der Charakteristiken konstant. Aus der ersten Gleichung erhalten wir zum einen die Identifizierung $\tau = t$, d.h. die Charakteristiken können durch die Zeit $t$ parametrisiert werden. Zum anderen die gewöhnliche Differentialgleichung $\dot{x} = \alpha$ mit der Lösung $x(\tau, \sigma) = g(\sigma) * \tau + s$ bzw. mit der Identifizierung dann $x(t) = g(\sigma) * t + \sigma$. Setzen wir die Definition von $g$ ein, so erhalten wir
	\begin{equation*}
		x(t) = \begin{cases}
			t + \sigma       & \sigma \le 0 \\
			(1-\sigma) t + \sigma & 0 \le \sigma \le 1 \\
			\sigma           & \sigma \ge 1
		\end{cases}
	\end{equation*}
	Für $t \le 1$ können wir jeden Fall umstellen und erhalten
	\begin{equation*}
		\sigma = \begin{cases}
			x - t & x - t \le 0\\
			\frac{x-t}{1-t} & x-t \ge 0 \land x \le 1 \\
			x & x \ge 1
		\end{cases}
	\end{equation*}
	Nach Konstruktion erhalten wir dann die Lösung
	\begin{equation*}
		u(x,t) = \alpha(\tau,\sigma) = g(\sigma(t,x)) = \begin{cases}
			1                   & x - t \le 0 \\
			1 - \frac{x-t}{1-t} & x-t \ge 0 \land x \le 1 \\
			0                   & x \ge 1
		\end{cases} 
	\end{equation*}
	Sei nun $t \in (0,1)$. Wir prüfen die Gültigkeit der Differentialgleichung:
	\begin{itemize}
		\item Sei $x - t \le 0$.  Dann ist also $u(x,t) = 1$ und $u_x = u_t \equiv 0$ und die Differentialgleichung erfüllt. 
		\item Ist $x-t \ge 0$ und $x \le 1$, dann ist $u(x,t) = 1 - \frac{x-t}{1-t}$ und $u_x(x,t) = \frac{1}{1-t}$ sowie $u_t(x,t) = \frac{(t-1)+(x-t)}{(1-t)^2}$. In die Differentialgleichung eingesetzt ergibt dies
		\begin{equation*}
			u_t + u u_x = \frac{(t-1)+(x-t)}{(1-t)^2} + \frac{1}{1-t} - \frac{x-t}{(1-t)^2} = 0 \quad \checkmark
		\end{equation*}
		\item Sei $x \ge 1$. Dann ist $u = u_x = u_t \equiv 0$ und die Differentialgleichung damit erfüllt.
	\end{itemize}
	Somit ist $u$ Lösung der Differentialgleichung für alle $x \in \R$ und $t \in (0,1)$.
\end{enumerate}



\begin{task}
	Betrachten Sie das Beispiel \enquote{Ampel (von rot auf grün)} aus der Vorlesung, modelliert durch
	\begin{equation*}
		\begin{aligned}
			u_t + F(u)_x &= 0 \qquad \text{für } (x,t) \in \R \times (0,\infty) \\
			u(x,0) &= g(x) \qquad \text{für alle } x \in \R
		\end{aligned}
	\end{equation*}
	mit 
	\begin{equation*}
		F(u) = u(60 - \frac{2}{5} u)\quad \text{und} \quad g(x) = \begin{cases} 150 & \text{für } x<0 \\ 0 & \text{für } x \ge 0\end{cases}
	\end{equation*}
	\begin{enumerate}
		\item In der Vorlesung wurde in Teil a) eine Funktion $u_a$ explizit gegeben. In Teil b) wurde eine weitere Lösung $u_b$ skizziert. Ermitteln Sie $u_b$ explizit.
		
		\emph{Hinweis}: Die charakteristischen Gleichungen geben $u_b$ auf eine großen Menge vor. Ermitteln Sie $u_b$ für die restlichen $(x,t)$, indem sie eine differenzierbare reelle Funktion $v$ derat bestimmen, dass $u_b(x,t) = v(\frac{x}{t})$ die Differentialgleichung löst.
		
		\item Überprüfen Sie für $u_a$ und $u_b$ entlang aller Sprungkurven die Rankine-Hugoniot-Bedingung und die Entropiebedingung.
	\end{enumerate}
\end{task}

\begin{enumerate}[label=(zu \alph*)]
	\item Für die skalare Erhaltungsgleichung erhalten wir die aus der Vorlesung bekannten Bezeichnungen mit $y = (x,t)$, $a(u,y) = (60 - \frac{4}{5} u \ , \ 1)$ und $b \equiv 0$. Außerdem sind die Startwerte der charakteristischen Gleichungen gegeben durch $x_0(s) = s$, $t_0(s) = 0$, und als Randwerte $g(s) = 150 * \one_{\menge{s < 0}}(s)$. Wir erhalten die charakteristischen Gleichungen
	\begin{equation*}
		\begin{aligned}
			\dot{y}(\tau, \sigma) &= \begin{pmatrix} \dot{x} \\ \dot{t} \end{pmatrix} = a(\alpha, y) = \begin{pmatrix} 60 - \frac{4}{5} u \\ 1 \end{pmatrix} \\
			\dot{\alpha}(\tau, \sigma) \overset{\text{DGl}}&{=} - b(\alpha, y) = 0  
		\end{aligned}
	\end{equation*}
	Aus der letzten Zeile der ersten Gleichung erhalten wir dabei die Identifizierung $t(\tau, \sigma) = \tau$. Die zweite charakteristische Gleichung löst sich unter Nutzung des Anfangswertes $\alpha(0,\sigma) = g(\sigma)$ zu $\alpha(\tau, \sigma) = g(\sigma)$. Setzen wir dies in die erste Zeile der ersten Gleichung ein, so erhalten wir $\dot{x}(\tau, \sigma) = 60 - \frac{4}{5} g(\sigma)$. Dies löst sich mit dem Anfangswert $x(0,\sigma) = \sigma$ zu 
	\begin{equation*}
		x(\tau, \sigma) = 60\tau - \frac{4}{5} g(\sigma) * \tau + \sigma = 
		\begin{cases}
			-60 t + \sigma & \sigma < 0 \\
			 60 t + \sigma & \sigma \ge 0
		\end{cases}
	\end{equation*}
	Für $\tau =  t > 0$ lässt sich dies umstellen zu
	\begin{equation*}
		\sigma(x,t) = 
		\begin{cases}
			x + 60 t & x < 0 \\
			x - 60 t & x \ge 0
		\end{cases}
	\end{equation*}
	Damit erhalten wir die Lösung 
	\begin{equation*}
		u(x,t) = \alpha(t, \sigma) = 
		\begin{cases}
			150 & \sigma(x,t) < 0 \\
			0   & \sigma(x,t) \ge 0
		\end{cases}
	\end{equation*}
	Da $t \in (0, \infty)$ ist, ist somit $u$ auf der Menge $\menge{(x,t) : t \in (0,\infty), x \notin (-60t, 60t)}$ eindeutig vorgegeben. Wir erweitern nun auf $t \in \R$ und betrachte $x \in (-60t, 60t)$. Dabei soll eine Funktion $v$ mit $u(x,t) = v(\frac{x}{t})$ bestimmt werden, die die PDE erfüllt. Die Funktion $u$ muss gemäß Vorlesung ihr Maximum auf der $t$-Achse, also für $x=0$ annehmen. Dementsprechend gilt $u(0,t) = v(\frac{0}{t}) = v(0) = u_\text{opt} = 75$. Da wir $x \in (-60t, 60t)$ betrachten, gilt $\frac{x}{t} = (-60,60)$. Die Randwerte sollten dabei mit minimaler (rechts) bzw. maximaler (links) Dichte gegeben sein, d.h. $v(-60) = 150$ und $v(60) = 0$. Diese drei Punkte definieren uns eine eindeutige Polynomfunktion zweiten Grades:
	\begin{equation*}
		v(\xi) = -\frac{5}{4s} \xi + 75 \qquad \text{auf } (-60, 60)
	\end{equation*}
	Damit können wir $u$ definieren als
	\begin{equation*}
		u(x,t) = v(\frac{x}{t}) = -\frac{5}{4} \frac{x}{t} + 75 \qquad \text{für } t \in \R \und x \in (-60t, 60t)
	\end{equation*}
	Damit gilt $u_x(x,t) = -\frac{5}{4} \frac{1}{t}$ und $u_t(x,t) = \frac{5}{4} \frac{x}{t^2}$. In die PDE eingesetzt liefert dies
	\begin{equation*}
		\begin{aligned}
			u_t + F(u)_x &= u_t + \brackets{60 - \frac{4}{5} u} u_x \\
			&= \frac{5}{4} \frac{x}{t^2} + \brackets{60 + \frac{x}{t} - 60} \brackets{-\frac{5}{4} \frac{1}{t}} 
			= \frac{5}{4} \frac{x}{t^2} -\frac{5}{4} \frac{x}{t^2} 
			= 0 \qquad \checkmark
		\end{aligned}
	\end{equation*} 
	Somit können wir $u$ schlussendlich definieren als
	\begin{equation*}
		u(t,x) = 
		\begin{cases}
			150 & x < -60t \\
			-\frac{5}{4} \frac{x}{t} + 75 & -60t \le x \le 60t \\
			0 & x > 60t
		\end{cases}
	\end{equation*}
	\item Definiere $u_b \defeq u$ von oben. Man sieht leicht, dass $u_b$ eine stetige Funktion ist. Daher existieren keine Sprungkurven und die Rankine-Hugoniot- bzw. Entropie-Bedingungen müssen nicht betrachtet werden. (Die Rankine-Hugoniot-Bedingungen werden trivialerweise trotzdem erfüllt, aber die Entropie-Bedingung nicht mehr).
	
	Aus der Vorlesung bekannt ist 
	\begin{equation*}
		u_a(x,t) = 
		\begin{cases}
			150 & x < 0 \\
			0   & x \ge 0
		\end{cases}
	\end{equation*}
	Wir betrachten die Sprungkurve $x=s(t)=0$. Dann ist auch $\dot{s} = 0$ und mit $F(u) = u \brackets{60 - \frac{2}{5} u}$ gilt
	\begin{equation*}
		\left.\begin{array}{rcl}
			\dsq{u} &=& 150 \\
			\dsq{F(u)} &=& 0
		\end{array}
		\right\} \follows \dsq{F(u)} = \dot{s} * \dsq{u}
	\end{equation*}
	Es ist $F'(u) = 60 - \frac{4}{5} u$ und damit $F'(u_l) = F'(150) = -60$ sowie $F'(u_r) = F'(0) = 60$. Somit ist die Entropie-Bedingung verletzt.
\end{enumerate}

\end{exercisePage}
\begin{exercisePage}
	
	\setcounter{taskcount}{11}
	
	\begin{task}
		Eine Funktion $u\in C^2(U)$ über einer offenen Menge $U\subset\Rn$ heißt harmonisch, falls $\Delta u = 0$ in $U$.
		\begin{enumerate}
			\item Zeigen Sie, dass die durch 
			\begin{equation*}
				u(x)= 
				\begin{cases}
					\ln(|x|) & \text{für } n=2, \\
					\frac{1}{|x|^{n-2}} & \text{für } n\geq 3
				\end{cases}
			\end{equation*}
			definierte Funktion auf $\Rn \setminus \menge{0}$ harmonisch ist.
			
			\item Es sei $\abb{u}{\Rn}{\R}$ harmonisch und $A\in \R^{n\times n}$ orthogonal. Zeigen Sie, dass durch $v(x)= u(Ax)$ eine harmonische Funktion $\abb{v}{\Rn}{\R}$ definiert wird.
		\end{enumerate}
	\end{task}

	\begin{enumerate}[label=(zu \alph*), leftmargin=*]
		\item Sei $n = 2$ und $x = (x_1, x_2) \in U$. Dann ist $u(x) = \ln \abs{x} = \ln(r(x))$ mit $r(x) \defeq \abs{x}$. Mit $r_{x_i}(x) = \frac{x_i}{\abs{x}}$ ist
		\begin{equation*}
			u_{x_i}(x) = \frac{1}{\abs{x}} * r_{x_i}(x) = \frac{x_i}{\abs{x}^2} = \frac{1}{x_1^2 + x_2^2}
		\end{equation*}
		Leiten wir weiter nach $x_i$ ab, dann erhalten wir
		\begin{equation*}
			u_{x_i x_i}(x) = \partial_{x_i} \frac{1}{x_1^2 + x_2^2} = \frac{x_1^2 + x_2^2 - x_i * 2x_i}{\abs{x}^4} = \frac{\abs{x}^2 - 2x_i^2}{\abs{x}^4}
		\end{equation*}
		und somit für den Laplace-Operator
		\begin{equation*}
			\Delta u(x) = x_{x_1 x_1} + u_{x_2 x_2} = \frac{\abs{x}^2 - 2x_2^2}{\abs{x}^4} + \frac{\abs{x}^2 - 2x_2^2}{\abs{x}^4}= 0 \quad \checkmark
		\end{equation*}
		Sei $n \ge 3$. Dann ist $u(x) = \abs{x}^{2-n}$. Es gilt
		\begin{equation*}
			\begin{aligned}
				u_{x_i}(x) &= (2-n) * \abs{x}^{1-n} * \frac{x_i}{\abs{x}} = (2-n) * \abs{x}^{-n} * x_i \\
				u_{x_i x_i} &= (2-n) * \partial_{x_i} \brackets{x_i * \abs{x}^{-n}} \\
				&= (2-n) * \brackets{\abs{x}^{-n} - n * \abs{x}^{-n-1} \frac{x_i}{\abs{x}} * x_i} \\
				&= (2-n) * \brackets{\abs{x}^{-n} - n * \abs{x}^{-n-2} * x_i^2} 
			\end{aligned}
		\end{equation*}
		Damit gilt für den Laplace-Operator
		\begin{equation*}
			\begin{aligned}
				\Delta u(x) = \sum_{i=1}^n u_{x_i x_i}(x) 
				&= (2-n) * \brackets{\sum_{i=1}^n \abs{x}^{-n} - n * \abs{x}^{-n-2} * \sum_{i=1}^n x_i^2} \\
				&= (2-n) \brackets{n * \abs{x}^{-n} - n* \abs{x}^{-n-2} * \abs{x}^2} \\
				&= 0 \qquad \checkmark
			\end{aligned}
		\end{equation*}
		
		\item Sei $\abb{u}{\Rn}{\R}$ harmonisch und $A \in \R^{n \times n}$ orthogonal, d.h. $\trans{A} = A^{-1}$. Betrachte
		\begin{equation*}
			\begin{aligned}
				v_{x_i}(x) &= Du(Ax) * \partial_i Ax = Du(Ax) * a_i = \sum_{k=1}^n u_{x_k}(Ax) * a_{ki} \\
				v_{x_i x_i}(x) &= \partial_{x_i} \brackets{\sum_{k=1}^n u_{x_k}(Ax) * a_{ki}} 
				= \sum_{k=1}^n \sum_{\ell=1}^n u_{x_k x_\ell}(Ax) * a_{ki} a_{\ell i}
			\end{aligned}
		\end{equation*}
		Aufgrund der Orthogonalität von $A$ ist $\sum_{i=1}^n a_{ki} a_{\ell i} = \delta_{k \ell}$. Somit ist
		\begin{equation*}
			\Delta v(x) = \sum_{i=1}^n v_{x_i x_i}(x) = \sum_{i=1}^n \sum_{k=1}^n \sum_{\ell=1}^n u_{x_k x_\ell}(Ax) * a_{ki} a_{\ell i}
			= \sum_{i=1}^n u_{x_i x_i}(Ax) 
			= \Delta u(Ax) = 0 \quad \checkmark
		\end{equation*}
		und daher auch $v$ harmonisch.
	\end{enumerate}
	
	\begin{task}
		Es sei $\Phi$ die Fundamentallösung der Laplace-Gleichung, d.\,h.
		\begin{equation*}
			\Phi(x) = \begin{cases} 
				-\frac{1}{2\pi} \ln (|x|) &\text{für } n=2,\\
				\frac{1}{n(n-2)\alpha(n)} \frac{1}{|x|^{n-2}} &\text{für } n \geq 3 
			\end{cases}
		\end{equation*}
		für $x\neq 0$. Zeigen Sie für $x\neq 0$ die Abschätzungen $\abs{D\Phi(x)} \le c \abs{x}^{1-n}$ und $\abs{D^2\Phi(x)} \le c \abs{x}^{-n}$.
	\end{task}
	
	Für $n=2$ gilt $\Phi_{x_i}(x) = -\frac{1}{2\pi} \frac{1}{\abs{x}} * \frac{x_i}{\abs{x}} = -\frac{1}{2\pi} * \frac{x_i}{\abs{x}^2}$. Leiten wir für $n \ge 3$ nach $x_i$ ab, so erhalten wir $\Phi_{x_i}(x) = \frac{2-n}{n(n-2) \alpha(n)} * \abs{x}^{1-n} \frac{x_i}{\abs{x}} = \frac{2-n}{n(n-2) \alpha(n)} * x_i * \abs{x}^{-n}$. Beide Fälle können wir mit einer universellen Konstante $c$ zusammenfassen und erhalten mit $\frac{x_i}{\abs{x}} \le 1$ auch folgende Abschätzung:
	\begin{equation*}
		\Phi_{x_i}(x) = c * x_i * \abs{x}^{-n} = c * \abs{x}^{1-n} * \frac{x_i}{\abs{x}} \le c * \abs{x}^{1-n} 
		\qquad (i = 1, \dots, n)
	\end{equation*}
	Damit erhalten wir für die erste Ungleichung
	\begin{equation*}
		\abs{D\Phi(x)} = \sqrt{\sum_{i=1}^n \Phi_{x_i}(x)^2} \le \sqrt{\sum_{i=1}^n \brackets{c * \abs{x}^{1-n}}^2}
		= c * n * \abs{x}^{1-n} = c * \abs{x}^{1-n}
	\end{equation*}
	Für die zweiten Ableitungen gilt dann entsprechend
	\begin{equation*}
		\begin{aligned}
			\Phi_{x_i x_j}(x) 
			&= c \brackets{\delta_{ij} \abs{x}^{-n} - n x_i \abs{x}^{-n-1} * \frac{x_j}{\abs{x}}} 
			= c * \abs{x}^{-n} \brackets{\delta_{ij} - n * x_i x_j * \abs{x}^{-2} } \\
			%
			\follows \quad
			%
			\abs{D^2 \Phi(x)} 
			&= \sqrt{\sum_{i=1}^n \sum_{j=1}^n \Phi_{x_i x_j}(x)^2}
			= \sqrt{\sum_{i=1}^n \sum_{j=1}^n \brackets{\abs{x}^{-n}}^2 * \brackets{c \delta_{ij} - n x_i x_j \abs{x}^{-2}}^2} \\
			&= \abs{x}^{-n} \sqrt{nc^2 - 2n \abs{x}^{-2} \sum_{i=1}^n x_i^2 + n^2 \abs{x}^{-4} \sum_{i=1}^n x_i^2 \sum_{j=1}^n x_j^2} \\
			&= \abs{x}^{-n} * \sqrt{nc^2 - 2n + n^2} \\
			&\le c \abs{x}^{-n}
		\end{aligned}
	\end{equation*}
\end{exercisePage}
\begin{exercisePage}
	\begin{task}
		Es sei $U \subset \Rn$ ein beschränktes Gebiet und $u \in C^2(U)\cap C(\overline U)$ eine Lösung von
		$$-\Delta u = f(u,x)\quad\text{in }U,\qquad u = g \quad\text{auf }\partial U,$$
		wobei $f$ und $g$ im Folgenden spezifisch gewählt werden.
		\begin{enumerate}
			\item Es sei $f(u,x) = u-u^3$ und $g=0$. Beweisen Sie \emph{mit elementaren Methoden} die Abschätzung $-1\leq u(x)\leq 1$ für alle $x\in \overline{U}$.
			
			\item Es sei $U = (-a,a)^n$ für ein $a>0$, $f(u,x) = -1$ und $g=0$. Finden Sie möglichst gute obere und untere Schranken für $u(0)$, indem Sie eine harmonische Funktion der Form $v=u+w$ betrachten, wobei $w$ geeignet zu wählen ist
			
			\item Es sei $U = B_1(0)$, $f(u,x) = h(x)$ mit $h,g \in C^1(\Rn,\R)$. Zeigen Sie, dass es eine von $n$, $h$, $g$ und $u$ unabhängige Konstante $c>0$ gibt, für die gilt:
			\begin{equation*}
				\max_{\quer{B_1(0)}} |u| \leq c * \left( \max_{\rand B_1(0)} |g| +\max_{\quer{B_1(0)}} |h| \right)
			\end{equation*}		
			\emph{Hinweis}: Betrachten sie $u-v$, wobei
			$v(x)=\max\limits_{\rand B_1(0)} |g| + \left(e^2 - e^{x_1+1} \right) \max\limits_{\quer{B_1(0)}} |h|$.
		\end{enumerate}
	\end{task}

	\begin{enumerate}[label=(zu \alph*), leftmargin=*]
		\item Es sei $f(u) = u - u^3$ und $g \equiv 0$. Da $U$ beschränktes Gebiet ist und $\quer{U}$ eine abgeschlossene Menge, ist also $\quer{U}$ kompakt. Da $u$ stetig auf $\quer{U}$ ist, nimmt $u$ ein Maximum in einem $x_0 \in \quer{U}$ an. Nehmen wir an es sei $u(x_0) > 1$. Es gilt $ \equiv g \equiv 0$ auf $\rand U$, d.h. um ein Maximum $u(x_0 > 1$ zu besitzen, muss $x_0 \in \inn U$ sein. Da $u(x_0)$ Maximum ist mit $u \in C^2(U)$, gilt $u_{x_i x_i}(x_0) \le 0$ für alle $i=1, \dots, n$. Nach Differentialgleichung ist dann also
		\begin{equation*}
			0 \le - \Delta u(x_0) =  - \sum_{i=1}^n u_{x_i x_i} (x_0) = u(x_0) - u^3(x_0) = u(x_0) \brackets{1 - u(x_0)^2} \overset{u(x_0) > 1}{<} 0
		\end{equation*}
		ein Widerspruch. Somit ist dann $u(x_0) \le 1$ und aufgrund der Maximalität von $u(x_0)$ auch $u(x) \le u(x_0) \le 1$ für alle $x \in \quer U$.
		
		Analog dazu nimmt $u$ auf $\quer U$ ein globales Minimum in $x_0 \in \quer U$ an, für welches nach gleicher Argumentation wie oben $x_0 \in \inn U$ gilt. Nehmen wir an, es sei $u(x_0) < -1$. Als Minimalstelle gilt $u_{x_i  x_i} (x_0) \ge 0$ für alle $i = 1, \dots, n$. Nach PDE gilt dann
		\begin{equation*}
			0 \ge - \Delta u(x_0) =  - \sum_{i=1}^n u_{x_i x_i} (x_0) = u(x_0) - u^3(x_0) = u(x_0) \brackets{1 - u(x_0)^2} \overset{u(x_0) < -1}{<} 0
		\end{equation*} 
		ein Widerspruch. Somit ist $u(x) \ge u(x_0) \ge -1$ für alle $x \in \quer U$ und schließlich gilt die Einschließung $-1 \le u(x) \le 1$ für alle $x \in \quer U$.
		
		\item Sei $a > 0$ und $U = (-a , a)^n$ ein $n$-dimensionaler (offener) Quader. Weiter sei $f \equiv -1$ und $g \equiv 0$. Gesucht sind \enquote{gute} (obere und untere) Schranken von $u(0)$.
		Wir betrachten eine harmonische Funktion $v$ der Form $v = u + w$, d.h. $0 = \Delta v = \Delta u + \Delta w = - f(u,x) + \Delta w = 1 + \Delta w$. Somit suchen wir nun eine Funktion $w$ mit $\Delta w = -1$. Eine Lösung dieser PDE erhalten wir beispielsweise mit $w(x) = - \frac{1}{2n} \abs{x}^2 = - \frac{1}{2n}\sum_{i=1}^n x_i^2$. 
		
		Wenden wir das Maximumsprinzip auf $v$ an, dann nimmt damit $v$ sein Maximum in einem $x_0 \in \rand U$ an. Aufgrund der Randwertbedingung gilt dort $u \equiv 0$. Wir betrachten oBdA den Randpunkt $x_0 = (a, 0, \dots, 0) \in \rand U$. Dieser minimiert $\sum_{i=1}^n x_i^2$, da jeder andere Randpunkt auch mindestens eine Koordinate $j \in \menge{1, \dots, n}$ mit $x_j = \pm a$ besitzt. Somit gilt dann schlussendlich
		\begin{equation*}
			u(0) = v(0) - w(0) = v(0) \le \max_{x \in U} v(x) \le \max_{x \in \rand U} \brackets{u(x) + w(x)} \le w(x_0) = - \frac{1}{2n} a^2
		\end{equation*}
		
		Analog liefert das Minimumsprinzip die Existenz des Minimum in $x_0 \in \rand U$. Nach Randwertbedingung gilt dort wieder $u(x_0) = 0$. Die Funktion $w$ wird auf dem Rand minimiert durch den Punkt $x_0 = (a, \dots, a) \in \rand U$ mit Minimalwert $w(x_0) = - \frac{1}{2n} \sum_{i=1}^n a^2 = - \frac{1}{2} a^2$. Analog zu oben gilt nun
		\begin{equation*}
			u(0) = v(0) - w(0) = v(0) \ge \min_{x \in U} v(x) \ge \min_{x \in \rand U} \brackets{u(x) + w(x)} \ge w(x_0) = - \frac{1}{2} a^2
		\end{equation*}
	\end{enumerate} 
	
	\begin{task}
		Für $t>0$ ist die Fundamentallösung der Wärmeleitungsgleichung auf $\Rn$ gegeben durch $u(x,t) = (c t)^{-n/2} e^{-\frac{|x|^2}{4t}}$ mit einer Konstante $c>0$. Zeigen Sie:
		\begin{enumerate}
			\item $u_t = \Delta u$ auf $\Rn \times \R_+$.
			\item $\lim_{t\to 0+} u(x,t)= 0$ für $x\neq 0$, $\lim_{t\to 0+} u(0,t) = \infty$ und $\lim_{t\to\infty} u(x,t)=0$.
			\item $\int_{\Rn} u(x,t)dx = 1$ für alle $t>0$ und $c= 4\pi$.
		\end{enumerate}
	\end{task}
\end{exercisePage}

\end{document}